\chapter{Conclusion}
\label{chap:conclusion}

In conclusion, while currently limited in their application in video games and at the same time challenging both users and developers working with them, BCIs hold promising potential for their future use in gaming. This is underlined by the fact that commercially available products that aid in the development of BCIs already exist and research in that field is encouraged by companies selling these products (such as \emph{g.tec} and \emph{EMOTIV}). The encouragement comes in the form of frequent hackathons, essentially short collaborative soft- and hardware development events, that are hosted by these companies and some of which even have an own category specifically for BCI gaming \cite{IEEESMC:hackathon, emotiv:hackathon}. % do we need to explain hackathons? (collaborative soft- and hardware development events)
With further development in both clinical and non-clinical BCIs and the possibility of new algorithms for calibrating BCIs being detected (especially in the field of neural networks) we might expect BCIs for gaming to be cheaper, have faster response times and thus finding a wider field of applications in a number of existing and upcoming games.